\section{Multi-subjects Time-varying graphical Lasso}

A network is composed of a group of nodes connected to each other by edges. 
Whether two nodes are connected or not depends on their dependence conditional on other nodes.
This is straightforward because an edge only reflects the relationship between two nodes regardless of the states of the rest. 
For Gaussian observations, a direct statistic showing this conditional dependence is the precision matrix, i.e. the inverse covariance matrices, 
whereby two nodes are conditionally independent to each other if and only if their corresponding precision matrix entries are zero\footnote{See Subsection \ref{subsec: pcorrandprec}}.
Based on this property, 
Hallac et al.\cite{hallac2017networkinferencetimevaryinggraphical} proposed Time-varying graphical lasso (TVGL).  
This is a method fully leveraging the time series observations i.i.d. from $\mathcal{N}(0,\Sigma(t))$ at different time stamps $t_1,\cdots,t_n$  
to acquire penalized likelihood estimator of precision matrices $\Theta=(\Theta_1, \cdots, \Theta_n)$
and hence to infer the snapshots of a time-varying network.
The method establishes the minimization goal by constructing the loss function consists of a negative log-likelihood with 
the $l_1$ spatial penalty on $\Theta$ and a convex spatial correlation penalty function on the difference between two adjacent precision matrices $\Theta_i$, $\Theta_{i+1}$. 
In addition, Hallac et al.\cite{hallac2017networkinferencetimevaryinggraphical} uses Alternating Direction Method of
Multipliers (ADMM) to derive an efficient algorithm to solve the optimization problem. 

TVGL infers the time-varying network from raw observations, 
whereas it is possible in real-life cases, for each time stamp, 
only sample covariance matrices respectively from different subjects are observable and the sample sizes of observations from each subject are unknown.
In abstract research backgrounds, we may deem a time series of positive definite real matrices as snapshots of a time-varying covariance matrix, 
with possibly multiple matrices for each time stamp. 
The typical research questions include inferring time varying cell-cell communication network from cell communication probability matrices, 
and discovering aging dynamic of brain structure connection. 
In this way, the observations consist of multiple sample covariance matrices but not single one per time stamp, 
on which TVGL can't perform directly.
In order to make TVGL applicable to the aforementioned cases, 
we adapt TVGL to a multi-subjects version called msTVGL and give the corresponding ADMM solver for $l_1$ and $l_2$ spatial penalty on the basis of the work of Hallac et al.. 
The adaption is equivalent to using the average of all observed covariance matrices for each time stamp. 
The msTVGL method is also applicative when in each time stamp the samples are collected from different subjects 
and can't be mix up directly as a whole considering heterogeneity between subjects. 
For example, when finding the aging dynamic of gene interaction, 
if two independent genes happen to express proactively or inactively at the same time in a subject, 
and cells collected from this subject are much more than the cells from the other subjects, 
the mixed data will contribute to a false high correlation between the two genes, 
thereby allowing TVGL to infer the existence of an unnecessary network edge.
However, the averaged within-subject covariance matrices used by msTVGL will not magnify the heterogeneity between subjects with unequal sample size, 
and hence reduce the strength of unreal network edges. 
Furthermore, the prominent difference that msTVGL distinguishes from TVGL is that 
msTVGL gives the partial correlation matrix at each time stamp to show the connection strength of nodes, 
and is able to show the change from $t_i$ to $t_{i+1}$
using the difference between the two partial correlation matrices. 

Below we first introduce the preliminary on the precision matrix and partial correlation, 
and then derive the updating algorithm of msTVGL with ADMM.